\FloatBarrier

\section{Conclusion and future work}

\FloatBarrier

\subsection{Conclusion}

The principal contribution of this work is a practical audit for \textbf{asymmetric training exposure between compared models}: a fault that resides in the comparison procedure rather than in any model's training pipeline, and that is therefore invisible to per-pipeline auditing. The hazard itself is documented in the foundation-model and benchmark-contamination literatures, but the standard leakage taxonomies are organised per pipeline and contain no category for it, while the existing remedies --- lineage tracing and contamination search --- presume documentation or corpus access that an external auditor rarely has.

Its three preconditions (Definition 2) are overlapping development corpora, group-structured data with within-group dependence, and an evaluation sample spanning at least one model's partition boundary. All three are common in appearance-based perception, because shared public corpora are the norm and identity is almost always the grouping key. Theorem 1 shows that under these conditions the reported score is a convex mixture of a memorisation and a generalisation regime whose weight is a sampling decision, and Corollary 1 shows the induced bias is signed and always favours the exposed model. In the system studied here the fault inflated one modality's apparent accuracy by 1.905$^{\circ}$, of which 0.425$^{\circ}$--1.068$^{\circ}$ (22 \%--56 \%) is attributable to training familiarity once the reference model prices the intrinsic difficulty of the affected stratum --- the interval spanning the two difficulty models of \EQREF{59}, which the available data cannot distinguish, and under both of which difficulty fails to explain the whole gap. This is a decomposition (Theorem 3) that lineage analysis cannot supply at all, since knowing that a model saw a corpus does not say how much of its advantage came from having seen it.

\textbf{Two instruments make the fault tractable in practice.} The per-group error ratio r(g) detects it from predictions alone --- the only assets an external audit reliably has are the two prediction streams and the group labels --- separating training membership at AUC 0.889 (p = 0.006). It adapts difficulty calibration, developed for membership inference, from example level to group level, which is what makes it stable enough to read off published predictions; Proposition 2 states the condition under which the intrinsic-difficulty term cancels exactly. It is a screening statistic rather than a membership oracle, and unlike black-box contamination proofs it carries no false-positive guarantee. A negative control on head pose, where neither model saw the evaluation groups, returns I$^{2}$ = 0.0 \% against 99.8 \% for the leaked gaze comparison, establishing that it stays silent when it should even under a fourfold swing in intrinsic difficulty. The six-step exposure-disjoint comparison protocol of \TABREF{5} then prevents the fault where the audit would otherwise have to catch it.

\textbf{Section 5.7 records what the correction was worth.} Under the uncorrected protocol the proposed gaze network scored level with its pretrained baseline and the decision was to keep it; under the corrected protocol the baseline is both more accurate and 7.1$\times$ smaller, and the decision reverses.

Evaluated under the corrected protocol, reproducible optimisation produced consistent gains, largest for voice (+9.81 pp) and head pose (--11.1 \% MAE); domain-targeted augmentation modelling the analytic train--serve shift \EQREF{34} reduced gaze angular error a further 9.5 \%; and a controlled scaling experiment established that the residual is bounded by \textbf{training identity count rather than sample count} --- the same group-structure property that makes the corpus leak-prone in the first place, by \EQREF{222}.

\textbf{The evaluation's principal analytical finding is that one mechanism explains the regression modalities' error structure.} Both networks operate in a shrinkage regime --- fitted gains of 0.48 on both gaze axes and 0.46--0.49 on the two narrow-support head-pose axes --- and that single parameter accounts for the isotropic radial growth of gaze error, the constant-ratio degradation of head-pose error, the axis-ranking inversions between absolute and variance-normalised metrics, and the absence of any localised failure mode. Inverting the contraction raises gaze angular error by 28.5 \%, so the networks sit near their minimum-MSE shrinkage point and improvement requires a better \emph{representation} rather than better \emph{calibration}. The classification modalities are limited differently: ROC-AUC of 0.9836 and 0.9961 against operating-point accuracies of 93.67 \% and 94.53 \% place their residual error in decision logic, which is why a confidence gate and a threshold move --- not additional capacity --- are what improve them.

\textbf{Reading those results through the deployed decision rules changes three conclusions} that the component metrics alone would suggest. Head pose survives an 8.74$^{\circ}$ yaw MAE because its error is unbiased and its gesture rule averages $\approx$20 frames inside an 18$^{\circ}$ hysteresis band, while its 0.822 gain manifests as a stable, learnable requirement to turn $\approx$36.5$^{\circ}$ to trigger a 30$^{\circ}$ threshold. Blink inverts the conventional cost asymmetry: because a click requires $\approx$13 consecutive frames above threshold, isolated false positives cannot fire one, so false negatives --- which break a deliberate closure run --- are the more damaging frame-level error, and thresholds below the default 0.50 are the defensible setting. Voice contributes a finding that is a design fault rather than a model fault: the weakest keyword in the vocabulary, \emph{go} at 0.877 precision, is bound to ``resume control'', the only command that reactivates a paused interface, and rebinding it to a phonetically distant token removes the largest error mode at zero computational cost.

\textbf{Measured against pretrained baselines on matched samples the proposed networks are less accurate, and this is reported without qualification:} 6.439$^{\circ}$ against 3.621$^{\circ}$ for gaze, 7.367$^{\circ}$ against 3.482$^{\circ}$ for head pose, 94.53 \% against 95.50 \% for voice. Their contribution is positional. For head pose and voice they hold competitive accuracy at 34$\times$ and 161$\times$ footprint reductions and an aggregate 6.68 ms per frame across three vision modalities, which is what permits four modalities to run concurrently in real time on commodity hardware without an accelerator. For gaze the pretrained alternative dominates on both accuracy and size, error correlation and oracle analysis show fusion cannot close the gap, and its adoption is recommended in place of the proposed network. Throughout, the attribution analysis establishes that the deployed networks compute the anatomically and acoustically appropriate operators, so the measured accuracies rest on identified transferable mechanisms rather than on dataset-specific shortcuts --- which is the property a system intended for a user it has never seen must actually possess.

\FloatBarrier

\subsection{Future work}

Six directions follow directly from the limitations of Section 6, ordered by how much each would strengthen the present conclusions.

\textbf{Replication of the detector across corpora and modalities.} GERA currently has one positive case and one negative control. The necessary next step is a set of comparisons spanning several corpora, several grouping keys, and --- critically --- cases in which leakage is known to be present but \emph{mild}, which would constrain the false-negative direction that the present negative control cannot. A synthetic protocol in which the exposed fraction $\Pi$ of \EQREF{38} is varied continuously from 0 to 1, with everything else held fixed, would produce exactly the calibration curve r($\Pi$) that the current study lacks and would permit a principled threshold on r to be proposed.

\textbf{Full leave-one-person-out gaze evaluation.} The single group-disjoint split of Section 6.1 should be replaced by leave-one-person-out over all fifteen MPIIGaze participants, which would convert the point estimates of Section 5.2 into fold-distributed estimates with measurable variance, and would supply fifteen rather than two observations of the identity-conditioned gain s(g) of \EQREF{220}. Since Section 5.2.4 identifies identity count as the binding constraint, this is also the experiment that would test that identification most directly.

\textbf{End-to-end measurement of the composite.} Section 6.5 is the largest gap between what is measured and what the user experiences. The three specific quantities required are: cursor error after per-user calibration and stabilisation, measured against known screen targets; click precision and recall in live sessions, which would also supply the AR(1) persistence $\varrho$ of \EQREF{123} that both the blink and head-pose arguments depend on; and gesture false-trigger rate under naturalistic head motion. All three are session-level measurements the deployed system already has the instrumentation to collect.

\textbf{Resolving the unexplained voice movement.} Section 5.1.1 records a 7.24 pp change in the voice figure between the published and fair-baseline columns despite an unchanged partition. Because the official Speech Commands split is already speaker-disjoint, this movement cannot be a leakage correction, and the present experiments do not identify its cause. Re-running the pre-correction voice configuration with the evaluation split logged explicitly would settle it in a single run, and until it is settled voice cannot serve as a second internal control alongside head pose --- which is the reason the control in Section 5.1.1 rests on one modality rather than two.

\textbf{Multi-seed evaluation and variance reporting.} Replacing the single-seed protocol of Section 6.3 with a small ensemble of seeds per configuration would attach a variance to every reported score and would settle whether the sub-1-pp differences currently declared inconclusive are real.

\textbf{The vocabulary redesign the analysis prescribes.} Section 5.4 identifies a zero-cost intervention --- rebinding ``resume control'' from \emph{go} to a phonetically distant token with fricative or stop-cluster onset --- predicted by \EQREF{263} to remove at least 50 \% of the total voice error mass. This prediction is testable at the cost of one evaluation run and would validate the entire decision-rule-aware framing, since it is a conclusion no component metric would have produced.

\textbf{Extending the exposure audit beyond pairwise comparison.} The present formulation compares two models. Realistic leaderboards compare many, with heterogeneous and undocumented exposure sets, and the natural generalisation is a matrix of pairwise ratios $r_{\mathrm{MM}}$'(g) whose structure would identify \emph{which} models in a set are exposed, not merely whether a given pair's comparison is contaminated. Whether the difficulty-cancellation of Proposition 2 survives that generalisation --- in particular whether a consistent reference can be identified from the matrix itself rather than assumed --- is an open theoretical question and, in our view, the most interesting one this work leaves behind.
